\subsection{行列式线丛与广义 \texorpdfstring{$\Theta$}{Θ}-除子}\label{src:9.4.3-subsection}
考虑图
\[
\begin{tikzcd}[ampersand replacement=\&,column sep=large]
V\arrow[d,dash]\&\mathcal F_{\mathrm{univ}}\arrow[d,dash]\&\mathcal L_V\arrow[dd,dash]\\
C\&C\times\cBun_{r,d}\arrow[l,"p_1"']\arrow[dr,"p_2"]\&{}\\
\&{}\&\cBun_{r,d},
\end{tikzcd}
\]
其中 $\mathcal F_{\mathrm{univ}}$ 是 $C\times\cBun_{r,d}$ 上的万有向量丛, $V$ 是 $C$ 上向量丛, $\mathcal L_V$ 是下述定义的 $\cBun_{r,d}$ 上的行列式线丛.
\begin{definition}[行列式线丛]\label{src:9.4.15}
关于 $V$ 的\newterm{行列式线丛}{determinantal line bundle} 定义为 $\cBun_{r,d}$ 上的线丛
\[
\mathcal L_V:=\det\left(\mathrm Rp_{2,*}(\mathcal F_{\mathrm{univ}}\otimes p_1^*V)\right)^\vee.
\]
\end{definition}
\begin{remark}[\texorpdfstring{$\mathcal L_V$}{LV} 良定义]\label{src:9.4.16}
因为投影 $p_2:C\times\cBun_{r,d}\to\cBun_{r,d}$ 可表, 射影且光滑, 相对维数为 $1$, 由上同调与基变换 (\bookref{A.6.2}), 对 $C\times\cBun_{r,d}$ 上每个向量丛 $E$, 导出推前 $\mathrm Rp_{2,*}E$ 是振幅位于 $[0,1]$ 的完美复形. 因而 $\mathrm R^ip_{2,*}E$ 在 $\cBun_{r,d}$ 上光滑局部可计算为向量丛二项复形 $[K^0\to K^1]$ 的上同调. 事实上, 它可整体计算为 $\cBun_{r,d}$ 上向量丛二项复形 $[K^0\to K^1]$ 的上同调; 见 \cite[命题 2.1.10]{HL10}. 无论整体还是光滑局部进行,
\[
\det\mathrm Rp_{2,*}E:=\det(K_0)\otimes\det(K_1)^\vee
\]
都是 $\cBun_{r,d}$ 上良定义的线丛. 关于复形行列式的更一般理论, 另见 \cite{KM76} 和 \cite[Tag 0FJI]{SP}.

取 $E=\mathcal F_{\mathrm{univ}}\otimes p_1^*V$, 可知行列式线丛 $\mathcal L_V$ 良定义. 对 $[F]\in\cBun_{r,d}(\Bbbk)$, 限制 $\mathcal L_V|_{[F]}$ 与 $\det\mathrm H^0(C,F\otimes V)^\vee\otimes\det\mathrm H^1(C,F\otimes V)$ 等同.
\end{remark}
\begin{exercise}[\texorpdfstring{$\mathcal L_V$}{LV} 的性质]\label{src:9.4.17}
\begin{enumerate}[(1)]
\item 证明存在从 $C$ 的 Grothendieck 群出发的群同态
\[
K_0(C)\to\operatorname{Pic}(\cBun_{r,d}),\qquad V\mapsto\mathcal L_V.
\]
特别地, 线丛 $\mathcal L_V$ 的同构类只依赖于 $\operatorname{rk}V$ 和 $\det V$.
\item 若 $V,W$ 是 $C$ 上秩与次数相同的向量丛, 则存在 $\underline{\mathcal P\!ic}^d$ 上的线丛 $\mathcal N$, 使得
\[
\mathcal L_W\cong\mathcal L_V\otimes\det^*\mathcal N,
\]
其中 $\det:\cBun_{r,d}\to\underline{\mathcal P\!ic}^d$ 是将 $F$ 映到 $\det F$ 的映射. 特别地, 对 $\xi\in\mathrm{Pic}(C)$, 行列式线丛 $L_V$ 到 $\cBun_{r,\xi}$ 的限制只依赖于 $V$ 的秩与次数.
\end{enumerate}
(习题~\bookref{9.4.42} 给出 $\mathcal N$ 的显式公式.)
\end{exercise}
\begin{remark}[秩 \texorpdfstring{$r=1$}{r=1} 的情形]\label{src:9.4.18}
取 $r=1$, 得到映射 $\mathrm{Pic}(C)\to\operatorname{Pic}(\underline{\mathcal P\!ic}^d)$, 将线丛 $M$ 映到 $\underline{\mathcal P\!ic}^d$ 上的行列式线丛 $\mathcal P_M$.
\end{remark}
\begin{remark}[\texorpdfstring{$V$}{V} 的数值条件]\label{src:9.4.19}
我们常要求对所有 $[F]\in\cBun_{r,d}(\Bbbk)$, 有 $\chi(C,F\otimes V)=0$, 因为这既使我们能构造 $\mathcal L_V$ 的整体截面, 又使它下降到良模空间. 由 Riemann--Roch 定理 (\bookref{9.1.2}), 此条件等价于
\begin{equation}\label{src:9.4.20}
d\operatorname{rk}(V)+r\deg(V)+r\operatorname{rk}(V)(1-g)=0,
\end{equation}
等价地, $\mu(F)+\mu(V)=g-1$.

令 $r'=r/\gcd(r,d)$, $d'=d/\gcd(r,d)$. 若向量丛 $V$ 对某整数 $m\geq1$ 满足
\[
\operatorname{rk}V=mr'\qquad\text{且}\qquad\deg V=mr'(g-1)-md',
\]
则 $V$ 满足 \bookeqref{9.4.20}.
\end{remark}
\textbf{$\mathcal L_V$ 的截面.} 设 $V$ 是 $C$ 上满足数值条件 \bookeqref{9.4.20} 的向量丛, $[K^0\to K^1]$ 是 $\cBun_{r,d}$ 上计算 $\mathrm Rp_{2,*}(\mathcal F_{\mathrm{univ}}\otimes p_1^*V)$ 的向量丛二项复形. 因为对所有 $[F]\in\cBun_{r,d}(\Bbbk)$, 有 $\chi(C,F\otimes V)=0$, $K^0,K^1$ 的秩相等! 所以
\begin{equation}\label{src:9.4.21}
s_V:\mathcal O_{\cBun_{r,d}}\to\det(K^0)^\vee\otimes\det(K^1)=:\mathcal L_V
\end{equation}
定义 $\mathcal L_V$ 的截面, 有时称为\newterm{$\Theta$-函数}{$\Theta$-function}. 虽然限制 $\mathcal L_V|_{\cBun_{r,\xi}}$ 只依赖于 $\operatorname{rk}V$ 和 $\deg V$ (习题~\bookref{9.4.17}(2)), 截面的限制 $s_V|_{\cBun_{r,\xi}}\in\Gamma(\cBun_{r,\xi},\mathcal L_V|_{\cBun_{r,\xi}})$ 确实依赖于 $V$.
\begin{lemma}\label{src:9.4.22}
截面 $s_V$ 在 $[F]\in\cBun_{r,d}(\Bbbk)$ 处不为零, 当且仅当 $\mathrm H^0(C,F\otimes V)=\mathrm H^1(C,F\otimes V)=0$.
\end{lemma}
\begin{proof}
截面不为零当且仅当 $K^0\otimes\kappa([F])\to K^1\otimes\kappa([F])$ 是同构, 这等价于 $F\otimes V$ 的 $\mathrm H^0$ 与 $\mathrm H^1$ 消失.
\end{proof}
\begin{remark}\label{src:9.4.23}
注意在数值条件 $\chi(C,F\otimes V)=0$ 下, $\mathrm H^0(C,F\otimes V)$ 消失等价于 $\mathrm H^1(C,F\otimes V)$ 消失. 上述引理将在 \S\bookref{9.6} 中 Faltings 的射影性证明中应用, 说明对适当的 $V$, $\mathcal L_V$ 无基点. 事实上, 我们将建立: 向量丛 $F$ 半稳定, 当且仅当存在向量丛 $V$, 满足 $\mathrm H^0(C,F\otimes V)=\mathrm H^1(C,F\otimes V)=0$; 这称为 Faltings 半稳定性判据 (\bookref{9.6.4}).
\end{remark}
\textbf{广义 $\Theta$-除子.}
\begin{definition}[广义 \texorpdfstring{$\Theta$}{Θ}-除子]\label{src:9.4.24}
设 $V$ 为向量丛, 对所有 $[F]\in\cBun_{r,d}(\Bbbk)$, 有 $\chi(C,F\otimes V)=0$, 即 \bookeqref{9.4.20} 成立. 关于 $V$ 的\newterm{广义 $\Theta$-除子}{generalized $\Theta$-divisor} 为
\[
\Theta_V:=\{[F]\in\cBun_{r,d}\mid\mathrm H^0(C,F\otimes V)\ne0\}.
\]
\end{definition}
\begin{proposition}[\texorpdfstring{$\Theta_V$}{ThetaV} 是除子]\label{src:9.4.25}
设 $r>0$ 和 $d$ 为整数, 令 $r'=r/\gcd(r,d)$, $d'=d/\gcd(r,d)$. 对 $m>0$ 以及一般的\footnote{``一般''是指存在稠密开子叠 $\mathcal V\subseteq\cBun_{mr',mr'(g-1)-md'}$, 使条件对每个 $V\in\mathcal V(\Bbbk)$ 成立.}秩为 $mr'$, 次数为 $mr'(g-1)-md'$ 的向量丛 $V$, $\Theta_V\subsetneq\cBun_{r,d}$ 是有效 Cartier 除子 (即余维为 $1$, 理想层为线丛的闭子叠), 且 $\mathcal L_V\cong\mathcal O_{\cBun_{r,d}}(\Theta_V)$.
\end{proposition}
\begin{remark}\label{src:9.4.26}
注意 $V$ 的数值条件强制 $\chi(C,F\otimes V)=0$ (见注记~\bookref{9.4.19}), 而 Faltings 半稳定性判据 (\bookref{9.6.4}) 蕴含每个这样的 $F,V$ 都半稳定.
\end{remark}
\begin{proof}
轨迹 $\Theta_V\subseteq\cBun_{r,d}$ 与截面 $s_V$ 的零轨迹等同 (引理~\bookref{9.4.22}), 由此赋予 $\Theta_V$ 概形结构. 为说明 $s_V\ne0$, 只需给出秩为 $r$, 次数为 $d$ 的向量丛 $F$, 满足 $\mathrm H^0(C,F\otimes V)=0$; 由引理~\bookref{9.4.22}, 这蕴含 $s_V([F])\ne0$. 由半连续性, 进一步只需找到一对具有指定数值不变量的向量丛 $V,F$, 满足 $\mathrm H^0(C,F\otimes V)=0$, 这正是习题~\bookref{9.4.27} 的结论. 另见 \cite{Hir88,Ghi84}. 当 $m\gg0$ 时, 此命题也来自独立证明的 Faltings 半稳定性判据 (\bookref{9.6.4}): 对每个 $F\in\cBun_{r,d}^{\mathrm{ss}}(\Bbbk)$, 存在向量丛 $V$, 使得 (等价地, 一般的 $V$ 满足) $\mathrm H^0(C,F\otimes V)=0$.
\end{proof}
\begin{exercise}[困难]\label{src:9.4.27}
设 $r>0$ 和 $d$ 为整数, 令 $r'=r/\gcd(r,d)$, $d'=d/\gcd(r,d)$. 证明存在秩 $r$, 次数 $d$ 的向量丛 $F$, 以及秩 $r'$, 次数 $r'(g-1)-d'$ 的向量丛 $V$, 使得 $\mathrm H^0(C,F\otimes V)=0$.
\end{exercise}
\begin{remark}[经典 \texorpdfstring{$\Theta$}{Θ}-除子]\label{src:9.4.28}
设 $L$ 为 $C$ 上次数为 $g-1-d$ 的线丛; 这保证对所有 $M\in\underline{\mathrm{Pic}}^d(C)$, 有 $\chi(C,M\otimes L)=0$. $\Theta$-除子定义为 $\Theta_L:=\{M\mid\mathrm H^0(C,M\otimes L)\ne0\}$, 可视为 Jacobian $\operatorname{Pic}^0(C)$ 上主极化 $\Theta$ 的平移.
\end{remark}
\textbf{$\mathcal L_V$ 的下降.}
\begin{proposition}\label{src:9.4.29}
设 $F$ 为秩 $r$, 次数 $d$ 的多重稳定向量丛. 若 $C$ 上向量丛 $V$ 满足 \bookeqref{9.4.20} 的数值条件 $d\operatorname{rk}(V)+r\deg(V)+r\operatorname{rk}(V)(1-g)=0$, 则 $\underline{\operatorname{Aut}}(F)$ 在纤维 $\mathcal L_V\otimes\kappa([F])$ 上平凡作用. 因而 $\mathcal L_V|_{\cBun_{r,d}^{\mathrm{ss}}}$ 下降为 $M_{r,d}^{\mathrm{ps}}$ 上线丛 $L_V$.
\end{proposition}
\begin{proof}
将 $F\cong\bigoplus_{i=1}^lF_i^{\oplus n_i}$ 写为互异稳定丛的直和, 则 $\underline{\operatorname{Aut}}(F)\cong\prod_i\mathrm{GL}_{n_i}$. 元素 $(g_i)\in\operatorname{Aut}(F)$ 在纤维
\[
\begin{aligned}
\mathcal L_V\otimes\kappa([F])&\cong\det\mathrm R\Gamma(C,F\otimes V)\\
&\cong\det\mathrm H^0(C,F\otimes V)\otimes\det\mathrm H^1(C,F\otimes V)^\vee\\
&\cong\bigotimes_i\left(\det\mathrm H^0(C,F_i\otimes V)\otimes\det\mathrm H^1(C,F_i\otimes V)^\vee\right)^{\otimes n_i}
\end{aligned}
\]
上的作用为乘以 $\prod_i\det(g_i)^{\chi(C,F_i\otimes V)}$. 每个 $F_i$ 的斜率为 $d/r$, 且 $\chi(C,F\otimes V)=0$, 因而 $\chi(C,F_i\otimes V)=0$, 作用平凡. $\mathcal L_V|_{\cBun_{r,d}^{\mathrm{ss}}}$ 下降为 $M_{r,d}^{\mathrm{ps}}$ 上线丛这一事实来自命题~\bookref{7.5.36}.
\end{proof}
